% IFN680 Project 7 technical report.
% Generated by tools/build_report.py from tools/report_template.tex. Every number below is
% printed by main_report.ipynb, except the design settings and the prior-work findings that
% tools/verify.py lists.
\documentclass[10pt,a4paper,twocolumn]{article}

\usepackage[margin=1.5cm,top=0.9cm,bottom=0.8cm]{geometry}
\usepackage{lmodern}
\usepackage[T1]{fontenc}
\usepackage[utf8]{inputenc}
\usepackage{microtype}
\usepackage{graphicx}
\usepackage{booktabs}
\usepackage{amsmath}
\usepackage[font=small,labelfont=bf,skip=3pt]{caption}
\usepackage{titlesec}
% Loaded last, as hyperref requires; it only fills the PDF's title and author fields.
\usepackage[hidelinks,pdftitle={Project 7: Improving DDPM},
            pdfauthor={Karan Rooprai, Nhu Hieu Nguyen}]{hyperref}

\titlespacing*{\section}{0pt}{4pt}{0.5pt}
\titleformat{\section}{\normalfont\bfseries}{\thesection.}{0.5em}{}
\setlength{\parskip}{0pt}
\setlength{\parindent}{1em}
\setlength{\columnsep}{0.6cm}
\setlength{\textfloatsep}{4pt}
\setlength{\dbltextfloatsep}{5pt}
\setlength{\intextsep}{3pt}
\setlength{\floatsep}{3pt plus 1pt minus 1pt}
\setlength{\tabcolsep}{4pt}
\pagestyle{empty}
\linespread{0.97}
% Never split a hyphenated word across a column or page break.
\brokenpenalty=10000

\renewcommand{\topfraction}{0.94}
\renewcommand{\dbltopfraction}{0.94}
\renewcommand{\bottomfraction}{0.6}
\renewcommand{\textfraction}{0.06}
\renewcommand{\floatpagefraction}{0.75}

\begin{document}

% The banner sits inside the \twocolumn optional argument, which is the only way a page-wide
% figure can appear on page 1 in the article class.
\twocolumn[{%
\begin{center}
{\large\bfseries Project 7: Improving DDPM. Does a Cosine Schedule Let a Diffusion Model Sample in
Fewer Steps?}\\[2pt]
{\bfseries Group 4}
\end{center}
\vspace{-2pt}
Student 1: Karan Rooprai \hfill Student 1 ID: n12498122 \\
Student 2: Nhu Hieu Nguyen \hfill Student 2 ID: n12194778\par
\vspace{4pt}
% \centering rather than a center environment, whose \topsep would add space above the figure.
{\centering\includegraphics[width=0.80\textwidth]{figures/figure_1_grids.pdf}\par}
\vspace{2pt}
{\small\textbf{Figure 1:} Twelve samples per digit from the linear-schedule (left) and
cosine-schedule (right) models at seed 0 and 1000 steps; images in the same place start from the
\emph{same} noise. The classifier reads __GRID_ACC_LINEAR__ (left) and
__GRID_ACC_COSINE__ (right) as the intended digit.\par}
\vspace{6pt}
}]

% A page-wide table can only float to the top of a later page, so it is declared here, on page 1,
% to appear at the top of page 2.
\begin{table*}[t]
\centering
\footnotesize
\begin{tabular}{@{}r rrll rrll@{}}
\toprule
& \multicolumn{4}{c}{DDPM sampler} & \multicolumn{4}{c}{DDIM sampler} \\
\cmidrule(lr){2-5}\cmidrule(l){6-9}
$K$ & Linear & Cosine & $\Delta$ [95\% CI] & $\Delta$, seeds 1, 2 & Linear & Cosine &
$\Delta$ [95\% CI] & $\Delta$, seeds 1, 2 \\
\midrule
__TABLE_ROWS__
\bottomrule
\end{tabular}
\caption{Fr\'echet distance to real test digits (lower is better; real-data floor __FLOOR__) at
$K$ sampling steps, seed-0 models, __N_A__ samples each. $\Delta$ is cosine minus linear with a
paired bootstrap interval, then at seeds 1, 2. Bold: better by the fixed rule.}
\end{table*}

\section{Introduction}
A denoising diffusion probabilistic model (DDPM; Ho et al., 2020) destroys an image in $T$ small
Gaussian steps and trains a network to predict the noise added, so that sampling can run the
process backwards from pure noise. The \emph{noise schedule} $\beta_1, \dots, \beta_T$ sets how
fast the share of signal left, $\bar\alpha_t = \prod_{s \le t}(1 - \beta_s)$, falls. Ho et al.\
raised $\beta_t$ linearly from $10^{-4}$ to 0.02 over $T = 1000$ steps. Nichol and Dhariwal (2021)
observed that the end of this process is ``too noisy'' to contribute much to sample quality, and
proposed a cosine schedule that destroys information more slowly. Here __NOISY_SHARE_LINEAR__\% of
the linear schedule's steps leave less than 1\% of the signal, against __NOISY_SHARE_COSINE__\% for
the cosine schedule, so a sampler that visits only $K$ of the 1000 steps wastes fewer of them on
near-pure noise. Hence the hypothesis: \emph{``A cosine schedule allows a denoising diffusion
probabilistic model to generate high-quality samples using fewer time steps than a linear
schedule.''}

Nichol and Dhariwal found that a linear-schedule model
``does not get much worse'' when up to 20\% of the reverse process is skipped. With the simple
loss used here, their cosine schedule improved log-likelihood on CIFAR-10 (3.37 to 3.26 bits/dim)
but not FID (2.90 to 3.05), and their sampling in far fewer steps came from \emph{learning} the
reverse variances; with fixed variances, as here, quality fell much faster as steps were removed.
Song et al.\ (2021) showed that DDIM, a deterministic sampler for the same trained network, gives
good samples 10 to 50 times faster. Lin et al.\ (2024) showed that common schedules and step
choices leave signal at the last step, which harms sampling in few steps. None of them isolates
what the schedule alone does to few-step sampling, which is the question tested here.

\section{Methodology}
\textbf{Baseline.} Tutorial 9.3's conditional DDPM, unchanged: \texttt{UNet\_cond} (__N_UNET__
parameters; its two poolings take $20 \to 10 \to 5$, so $20\times20$ images fit), with $t/T$ and
the class each embedded as 32 numbers added at every level. It predicts the added noise (mean
squared error), trained with Adam ($2\times10^{-4}$) in batches of 128 for __EPOCHS__ epochs on
the __N_TRAIN__ training images scaled to $[-1, 1]$, with $T = 1000$.

\textbf{Cosine schedule.} The only change: $\bar\alpha_t = f(t)/f(0)$ with
$f(t) = \cos^2\!\big(\tfrac{t/T + s}{1 + s}\cdot\tfrac{\pi}{2}\big)$, $s = 0.008$, and
$\beta_t = 1 - \bar\alpha_t/\bar\alpha_{t-1}$ clipped at 0.999 (Nichol and Dhariwal, eq.~17).
Training is otherwise identical: for seeds 0, 1 and 2, both start from the same weights,
batches, steps and noise.

\textbf{Sampling in $K$ steps.} A $K$-step sampler visits $K$ evenly spaced steps rounded to
integers (Nichol and Dhariwal, \S4) and jumps from $t$ to the next visited $t'$ with
$\beta = 1 - \bar\alpha_t/\bar\alpha_{t'}$. Two samplers use it: the tutorial's stochastic step
(DDPM) and the deterministic DDIM step. Both form the clean-image estimate $\hat x_0$ from the
predicted noise and clip it to $[-1, 1]$, as Ho et al.'s released sampler does; unclipped at
$K = T$ the step equals the tutorial's to rounding error. The clip matters for the cosine schedule,
whose last $\beta$ is 0.999: the tutorial's form multiplies an error in the first noise estimate by
__GAIN1000_COSINE__ at $K = 1000$ and __GAIN10_COSINE__ at $K = 10$ (linear: __GAIN10_LINEAR__).

\section{Experiments}
\textbf{E1} (Figure~1): 12 samples per digit from each seed-0 model at $K = 1000$.
\textbf{E2} (Table~1, Figure~2): each seed-0 model at $K = 10, 20, 50, 100, 250$ with both
samplers and at 1000 with DDPM, __N_A__ samples per set, drawn under the labels of a
class-stratified test subset A from the same starting noise for both schedules. \emph{Quality} is
the Fr\'echet distance (FD; Heusel et al., 2017) to a disjoint real subset B with the same class
mix, in the 128 features of a separate classifier (__CLF_ACC__ test accuracy), as Salvado (2026,
pp.~42, 87) proposes; real A against B gives the floor, __FLOOR__. \emph{Class consistency} is that
classifier's accuracy on the intended labels, and \emph{variety} the within-class feature spread
relative to real digits. \textbf{E3}: seeds 1 and 2 at every $K < 1000$. \textbf{A rule fixed in
advance}: at a given $K$ and sampler, a schedule is better if its FD is lower at seed 0 with a
95\% paired bootstrap interval (1,000 resamples) excluding zero, and lower at seeds 1 and 2 too.
\textbf{E4}: seconds per epoch, and the FD of seed 0's snapshots at epochs 10, 25, 50, 100 and
__EPOCHS__ (DDIM, $K = 50$). \textbf{E5}: each model's noise-prediction error on test images at 20
fixed steps. \textbf{E6}, run after the verdict: the seed-0 DDPM runs with Ho et al.'s smaller
posterior variance $\tilde\beta$.

\section{Discussion}
__QUALITATIVE__

__VERDICT_DDPM__ __VERDICT_DDIM__ __FEWEST_SENTENCE__ __FULL_SENTENCE__

\setcounter{figure}{1}
\begin{figure}[t]
\centering
\includegraphics[width=\columnwidth]{figures/figure_2_steps.pdf}
\caption{FD (top) and class consistency (bottom) against $K$ at seed 0, with seeds 1 and 2 as
small dots; the dotted lines are real test digits.}
\end{figure}

__INTERPRETATION__

\textbf{Training time and image quality.} __TIME_SENTENCE__ __LOSS_SENTENCE__

__SYNTHESIS__

\section{Conclusion}
__CONCLUSION__

\section*{References}
\footnotesize
\setlength{\parindent}{0pt}
\newcommand{\bib}[1]{\par\hangindent=1em #1}
\bib{Heusel, M. et al. (2017). GANs trained by a two time-scale update rule converge to a local
Nash equilibrium. \emph{NeurIPS}.}
\bib{Ho, J., Jain, A. and Abbeel, P. (2020). Denoising diffusion probabilistic models.
\emph{NeurIPS}.}
\bib{Lin, S., Liu, B., Li, J. and Yang, X. (2024). Common diffusion noise schedules and sample
steps are flawed. \emph{WACV}.}
\bib{Nichol, A. and Dhariwal, P. (2021). Improved denoising diffusion probabilistic models.
\emph{ICML}.}
\bib{Salvado, O. (2026). Week 9: transforming distribution (GAN, diffusion, flow). IFN680
lecture, QUT.}
\bib{Song, J., Meng, C. and Ermon, S. (2021). Denoising diffusion implicit models. \emph{ICLR}.}

\end{document}
